\section{Introduction}
\label{sec:introduction}

Let $(M,g)$ be a closed connected smooth Riemannian manifold of dimension $d$.
We write $-\Delta_g u=\lambda u$ for a nonzero real Laplace eigenfunction
with $\lambda>0$, and $Z_u=\{x\in M:u(x)=0\}$ for its nodal set. Yau's
nodal conjecture in the smooth category predicts constants $c_g,C_g>0$,
depending on this one metric, such that
\begin{equation}\label{eq:intro-yau}
 c_g\sqrt\lambda\le\mathcal H_g^{d-1}(Z_u)
              \le C_g\sqrt\lambda
\end{equation}
for every such eigenfunction. The distinction between a fixed metric and a
metric chosen anew at each eigenvalue is essential. The lower and upper
inequalities are also logically separate. This paper disproves the upper
inequality in two low-dimensional smooth settings.

\begin{theorem}\label{thm:main}
Let $g_*$ be the round metric on $\Sph^3$. Every neighborhood of $g_*$ in
the $C^\infty$ topology contains a smooth Riemannian metric $g_\infty$ for
which there are nonzero smooth real functions $u_j$ and positive numbers
$\lambda_j$ satisfying
\[
 -\Delta_{g_\infty}u_j=\lambda_j u_j,\qquad
 \lambda_j\longrightarrow\infty,\qquad
 \frac{\HH^2_{g_\infty}(Z_{u_j})}{\sqrt{\lambda_j}}
       \longrightarrow\infty.
\]
\end{theorem}

\begin{theorem}\label{thm:four}
There are a smooth Riemannian metric $g_\infty$ on $S^2\times\T^2$,
nonzero smooth real functions $u_j$, and positive numbers
$\lambda_j\to\infty$ such that
\[
 -\Delta_{g_\infty}u_j=\lambda_j u_j,\qquad
 \frac{\cH^3_{g_\infty}(Z_{u_j})}{\sqrt{\lambda_j}}
       \longrightarrow\infty.
\]
\end{theorem}

Theorem~\ref{thm:main} gives a negative resolution of the upper-bound
part of the smooth nodal conjecture already in dimension three, with
metrics arbitrarily close to the round sphere. Theorem~\ref{thm:four}
provides a distinct four-dimensional topology and metric realization.
Neither theorem asserts a power-law excess, and neither contradicts the
lower-bound part of \eqref{eq:intro-yau}.
In dimensions at least five, a separate companion constructs fixed-metric
counterexamples with a fixed positive power-law excess
\cite[Theorem~1.1 and Corollary~5.6]{OpenAINodalPower2026}.

\paragraph{Context.}
Yau formulated the nodal measure problem for Laplace eigenfunctions on
closed manifolds \cite{Yau1982}. Donnelly and Fefferman established both
sharp estimates for real-analytic metrics
\cite[Theorem~1.2]{DonnellyFefferman1988}.
For smooth metrics, Hardt and Simon obtained a general non-polynomial
upper estimate \cite[Section~5]{HardtSimon1989}. Logunov later proved
a polynomial upper estimate in dimensions at least three
\cite{LogunovUpper2018}, as well as the sharp lower estimate
\cite[Theorem~1.1]{LogunovLower2018}. The gap between polynomial growth
and the conjectured square-root upper bound is where the examples here lie.

The smooth surface case has a separate quantitative history.
Br\"uning proved the sharp square-root lower estimate
\cite[Satz~2]{Bruning1978}. Donnelly and Fefferman obtained the upper
estimate $C_g\lambda^{3/4}$ \cite{DonnellyFefferman1990}, and Logunov
and Malinnikova improved its exponent to $3/4-\beta$ for a universal
$\beta\in(0,1/4)$ \cite{LogunovMalinnikova2018}.
Under intermediate regularity hypotheses, Hezari proved bounds for
Gevrey metrics and a specified quasianalytic class
\cite[preprint version, Theorems~1.1 and~1.2]{Hezari2023}; these are not statements for
arbitrary smooth metrics.

Metric flexibility has also been used to prescribe nodal geometry.
On a closed manifold of dimension at least three, Enciso and Peralta-Salas
realized any smooth connected closed orientable separating hypersurface
as the nodal set of a first nonconstant eigenfunction by choosing the metric
\cite[preprint version, Theorem~1.1]{EncisoPeraltaSalas2015}.
Their work with Steinerberger prescribes any smooth connected closed
separating hypersurface in the same dimensions, within each conformal class
and with fixed total volume, up to an ambient diffeomorphism arbitrarily
close to the identity in $C^0$
\cite[preprint version, Theorem~1.1]{EncisoPeraltaSalasSteinerberger2017}.
On a fixed round sphere or flat torus, high-energy eigenfunctions can
have nodal components of prescribed complicated topology
\cite{EncisoPeraltaSalasTorres2018}. Krauz has prescribed individual
nodal configurations on the sphere by metric choice \cite{Krauz2025}.
These results address nodal topology, including at high frequencies on
fixed analytic metrics. The distinction in Theorems~\ref{thm:main}
and~\ref{thm:four} is the unbounded normalized nodal measure along a
sequence of eigenfunctions for one fixed smooth metric.

\paragraph{Two constructions and their common logic.}
Each branch starts from a background metric with explicit high-frequency
eigenfunctions and a region left unchanged by all local corrections.
A smooth logarithmic envelope is corrugated in transverse directions so
that its integrated gradient becomes arbitrarily large without losing
the strict Hessian inequality needed for localized complex waves.
Finite phase and transport jets give quasimodes with arbitrarily small
weighted residual. This construction is related to complex-phase WKB
quasimodes \cite[Section~3]{DenckerSjostrandZworski2004}; the weighted
finite-order estimates required here are proved in the wave sections. A Gaussian superposition makes the real value and
first derivative jointly nonvanishing and supplies many strict
opposite-sign tests along short coordinate segments. These tests form a
lower semicontinuous score bounded by the nodal Hausdorff measure, even
when the zero set is not assumed regular. Random-wave studies of typical
nodal volume provide a related probabilistic backdrop
\cite{RudnickWigman2007,Wigman2008}; the argument here selects one wave
with a large stable sign score and proves its realization under a metric
change, rather than importing a random-wave nodal-volume formula.

The obstacle is to replace each quasimode by an exact eigenfunction of an
\emph{ordinary} Laplace--Beltrami operator while retaining every sign.
In the three-sphere branch, a scalar conductivity and independent density
first correct the equation. Division by a positive solution of a
semilinear elliptic equation makes them satisfy the three-dimensional
metric determinant identity in the active region; a transverse tensor
change repairs the remaining defect in a clean annulus. In the
$S^2\times\T^2$ branch, a conformal scalar correction has coercive
linearization $\Delta_h-2\lambda$, and a change tangent to the level
sets enforces the four-dimensional determinant identity. These are
separate exactification arguments, not a dimension-independent
realization lemma.

Figure~\ref{fig:two-repairs} isolates the two dimension-specific repairs.
The first converts a scalar conductivity and density into a metric; the
second uses a conformal correction before repairing a rank-three tangent
bundle. In both rows, the last tensor change annihilates the gradient of
the chosen eigenfunction and therefore leaves its corrected equation
unchanged.
\begin{figure}[htbp]
\centering
\begin{tikzpicture}[
  every node/.style={draw,align=center,font=\footnotesize,minimum height=10mm,text width=25mm},
  >=Latex,node distance=2.2mm]
  \node (a) {three-sphere\\quasimode};
  \node[right=of a] (b) {scalar weighted\\correction};
  \node[right=of b] (c) {positive factor\\and division};
  \node[right=of c] (d) {two-plane\\repair};
  \node[right=of d] (e) {exact metric\\eigenpair};
  \node[below=7mm of a] (f) {$S^2\times\T^2$\\quasimode};
  \node[right=of f] (g) {tensor weighted\\correction};
  \node[right=of g] (h) {conformal\\factor};
  \node[right=of h] (i) {three-plane\\repair};
  \node[right=of i] (j) {exact metric\\eigenpair};
  \foreach \x/\y in {a/b,b/c,c/d,d/e,f/g,g/h,h/i,i/j}
    \draw[->] (\x) -- (\y);
\end{tikzpicture}
\caption{Two routes from localized quasimodes to exact eigenfunctions of
ordinary metrics. The repairs preserve the strict sign tests used later
to pass to one fixed limiting metric.}
\label{fig:two-repairs}
\end{figure}

Finally, a perturbation preserving the selected eigenfunction makes its
eigenvalue simple. A strict lower bound for the sign score then persists for all metrics
in a neighborhood. Nested metric neighborhoods retain each previous
witness while inserting the next one, and converge to the single smooth
metric required by each theorem. The large-frequency thresholds may
depend on each fixed envelope; no uniform spectral gap is required.
Uhlenbeck's generic simplicity theorem \cite{Uhlenbeck1976} is context
for this spectral step, not a substitute for the selected-mode-preserving
perturbations proved here.

Taking the product of the three-sphere metric with any fixed closed
Riemannian manifold and pulling back its eigenfunctions gives further
qualitative failures in higher dimensions. This elementary extension
does not assert a power-law excess in dimensions three or four.

Together with the surface theorem \cite[Theorem~1.1]{OpenAINodalSurfaces2026},
Theorem~\ref{thm:main} and this product observation give the exact universal
dimension threshold: among dimensions $d\ge2$, the upper inequality in
\eqref{eq:intro-yau}, with a metric-dependent constant uniform over all
nonzero real eigenfunctions of positive eigenvalue, holds for every closed
connected smooth Riemannian manifold precisely when $d=2$. For each
$d\ge3$, there is one fixed smooth metric with a violating sequence.

\paragraph{Organization and conventions.}
The three-sphere proof first constructs amplifiable profiles in
Section~\ref{sec:profiles}, waves and signs in Section~\ref{sec:waves},
an exact ordinary metric in Section~\ref{sec:exactification}, and a
fixed-metric limit in Sections~\ref{sec:pinning}--\ref{sec:persistence}.
The $S^2\times\T^2$ proof then develops its own envelopes in
Section~\ref{sec:envelopes}, waves in Section~\ref{sec:four-waves},
sign score in Section~\ref{sec:signs}, metric correction in
Section~\ref{sec:correction}, and spectral limit in
Section~\ref{sec:spectral}. Throughout, gradients and divergences
carry the displayed metric unless stated otherwise. Hausdorff measure
is taken for the Riemannian distance, with a fixed standard
normalization. All finite derivative and residual orders are chosen
before sending the frequency to infinity.
